\documentclass[11pt]{article}

\usepackage[margin=1in]{geometry}
\usepackage{amsmath,amssymb,amsthm}
\usepackage{mathpartir}
\usepackage{array}

%% PFPL conventions: operators sans-serif, variables italic,
%% arguments separated by ';', binding by '.'
\newcommand{\op}[1]{\mathsf{#1}}
\newcommand{\Exp}{\op{Exp}}
\newcommand{\Typ}{\op{Typ}}

\theoremstyle{definition}
\newtheorem{definition}{Definition}

\pagestyle{empty}

\begin{document}

\begin{definition}[Abstractor]
An \emph{abstractor} of valence $s_1,\dots,s_k.s$ is a phrase
\[
  x_1,\dots,x_k.a ,
\]
abbreviated $\vec{x}.a$, where $x_1,\dots,x_k$ are distinct variables with
$x_i$ of sort $s_i$, and $a$ is an abt of sort $s$ in which each $x_i$ is
bound. For $k = 0$ it is just $a$.
\end{definition}
\bigskip
\noindent\textbf{How to write it.}

\begin{center}
\begin{tabular}{@{}>{\ttfamily\small}l@{\qquad}l@{}}
\multicolumn{1}{@{}l@{\qquad}}{\textbf{source}} & \textbf{output} \\[4pt]
\hline \\[-6pt]
x.a                                & $x.a$ \\[4pt]
\textbackslash vec\{x\}.a          & $\vec{x}.a$ \\[4pt]
x\_1,\textbackslash dots,x\_k.a    & $x_1,\dots,x_k.a$ \\[4pt]
\textbackslash Exp.\textbackslash Exp
                                   & $\Exp.\Exp$ \\[4pt]
(\textbackslash Exp,\ \textbackslash Exp.\textbackslash Exp)\textbackslash Exp
                                   & $(\Exp,\ \Exp.\Exp)\Exp$ \\[4pt]
\textbackslash op\{let\}(a\_1 \textbackslash,;\textbackslash, x.a\_2)
                                   & $\op{let}(a_1 \,;\, x.a_2)$ \\[4pt]
\textbackslash op\{lam\}[\textbackslash tau](x.a)
                                   & $\op{lam}[\tau](x.a)$ \\[4pt]
[b/x]a                             & $[b/x]a$ \\[4pt]
a =\_\textbackslash alpha b        & $a =_\alpha b$ \\
\end{tabular}
\end{center}

\bigskip
\noindent\textbf{Formation rule (for reference).}

\begin{mathpar}
\inferrule*[right=op]
  { \mathcal{X}, \vec{x}_1 \vdash a_1 : s_1 \\ \cdots \\
    \mathcal{X}, \vec{x}_n \vdash a_n : s_n }
  { \mathcal{X} \vdash o(\vec{x}_1.a_1 \,;\, \dots \,;\, \vec{x}_n.a_n) : s }
\end{mathpar}

\noindent
where $o$ has arity $(\vec{s}_1.s_1,\dots,\vec{s}_n.s_n)s$.

\end{document}
